\thispagestyle{otFront}
\phantomsection\label{references}\pdfbookmark[0]{Références}{bm-references}
\otPageTitre{Références}{France · Première générale · Spécialité physique-chimie}
\fontsize{10.5}{14.5}\selectfont
\otIntertitre{Le programme de référence}
Le livre s'appuie sur le programme de l'enseignement de spécialité de physique-chimie de la
classe de première de la voie générale, publié au \textbf{Bulletin officiel spécial n° 1 du
22 janvier 2019} (arrêté du 17 janvier 2019, NOR MENE1901635A) et en vigueur depuis la rentrée
2019. Aucun nouveau programme de physique-chimie de première n'a été publié depuis : le
Bulletin officiel n° 14 du 2 avril 2026 ne modifie que les programmes de mathématiques du lycée.
Les quatre thèmes du programme structurent les parties du livre ; la partie « Mesure et
incertitudes » fait l'objet du chapitre 21.\par\vspace{3mm}
\href{https://www.education.gouv.fr/bo/19/Special1/MENE1901635A.htm}{\textcolor{otNavy}{Programme officiel sur education.gouv.fr}}\par\vspace{5mm}
\otIntertitre{La suite en terminale}
Les notions de première sont reprises en terminale et peuvent être mobilisées lors de l'épreuve
écrite de spécialité. Une spécialité arrêtée en fin de première est évaluée en contrôle continu,
à partir des résultats de l'année.\par\vspace{5mm}
\otIntertitre{Les choix pédagogiques}
Les exercices, les défis et les devoirs surveillés sont des créations pour ce livre : ils
s'inspirent des devoirs surveillés, devoirs maison et contrôles donnés en première, sans en
reproduire aucun. Les défis exigeants sont facultatifs ; les outils qui dépasseraient le
programme sont fournis dans l'énoncé. Les données réelles sont accompagnées de leur source,
consultée à la date de l'édition.\par\vspace{3mm}
Les schémas, courbes et tableaux sont vectoriels. Les programmes Python peuvent être exécutés
séparément. Les résultats numériques sont donnés avec un nombre de chiffres significatifs
cohérent avec les données.\par\vspace{5mm}
{\fontsize{9}{12}\selectfont\color{otMuted}Sources officielles consultées le 23 octobre 2026.
La charte reprend le modèle OptiTAB : couleurs, typographie et repères de difficulté.}
\clearpage
